% Part: lambda-calculus
% Chapter: introduction
% Section: term-revisited

\documentclass[../../../include/open-logic-section]{subfiles}

\begin{document}

\olfileid[id]{lam}{syn}{tr}
\olsection{Suku sebagai Kelas Ekuivalensi-$\alpha$}

Mulai sekarang, kita akan memandang suku hingga ekuivalensi-$\alpha$. Artinya,
ketika kita menulis suatu suku, yang kita maksud adalah kelas
ekuivalensi-$\alpha$ tempat suku itu berada. Sebagai contoh, kita menulis
$\lambd[a][\lambd[b][a c]]$ untuk himpunan semua suku yang $\alpha$-ekuivalen
dengannya, seperti $\lambd[a][\lambd[b][a c]]$,
$\lambd[b][\lambd[a][b c]]$, dan seterusnya.

Selain itu, sementara pada bagian-bagian sebelumnya huruf seperti $N, Q$
digunakan untuk menyatakan suatu suku, mulai sekarang kita menggunakannya
untuk menyatakan suatu kelas, dan kelas-kelas inilah, alih-alih suku, yang
menjadi objek kajian kita selanjutnya. Huruf seperti $x, y$ tetap menyatakan
variabel.

Kita juga menggunakan notasi $\rep{M}$ untuk menyatakan sebarang
!!{element} dari kelas~$M$, serta $\rep{M}[0], \rep{M}[1]$, dan seterusnya,
jika kita memerlukan lebih dari satu.

Kita menggunakan kembali notasi suku untuk menyederhanakan perumusan. Kita
mempunyai definisi berikut pada kelas-kelas:
\begin{defn}
  \begin{enumerate}
  \item $\lambd[x][N]$ didefinisikan sebagai kelas yang memuat
    $\lambd[x][\rep{N}]$.
  \item $PQ$ didefinisikan sebagai kelas yang memuat $\rep{P}\rep{Q}$.
  \end{enumerate}
\end{defn}

Tidak sulit untuk melihat bahwa keduanya terdefinisi dengan baik, karena
konversi-$\alpha$ kompatibel.

\begin{defn} \ollabel{def:fv}
  \emph{Variabel bebas} dari suatu kelas ekuivalensi-$\alpha$~$M$, atau
  $\FV{M}$, didefinisikan sebagai $\FV{\rep{M}}$.
\end{defn}

Definisi ini terdefinisi dengan baik karena $\FV{\rep{M}[0]} =
\FV{\rep{M}[1]}$, sebagaimana ditunjukkan dalam \olref[alp]{thm:fv}.

Kita juga menggunakan kembali notasi substitusi pada kelas-kelas:
\begin{defn} \ollabel{def:sub}
  \emph{Substitusi} $R$ bagi $y$ dalam $M$, atau $\Subst{M}{R}{y}$,
  didefinisikan sebagai $\Subst{\rep{M}}{\rep{R}}{y}$, untuk sebarang
  $\rep{M}$ dan $\rep{R}$ yang membuat substitusi tersebut terdefinisi.
\end{defn}

Definisi ini juga terdefinisi dengan baik, sebagaimana ditunjukkan dalam
\olref[alp]{cor:sub}.

Perhatikan bagaimana definisi ini sangat menyederhanakan penalaran kita.
Sebagai contoh:
\begin{align}
  \Subst{\lambd[x][x]}{y}{x} & \ollabel{eq:1}\\
  &= \Subst{\lambd[z][z]}{y}{x} \ollabel{eq:2}\\
  &= \lambd[z][\Subst{z}{y}{x}] \\
  &= \lambd[z][z]
\end{align}

\olref{eq:1} tidak terdefinisi jika kita masih memandangnya sebagai
substitusi pada suku. Akan tetapi, sebagaimana disebutkan sebelumnya, sekarang
kita memandangnya sebagai substitusi pada kelas; itulah sebabnya
\olref{eq:2} dapat dilakukan: kita dapat mengganti $\lambd[x][x]$ dengan
$\lambd[z][z]$ karena keduanya termasuk dalam kelas yang sama.

Atas alasan yang sama, mulai sekarang kita akan mengandaikan bahwa wakil yang
kita pilih selalu memenuhi syarat yang diperlukan bagi substitusi. Sebagai
contoh, ketika melihat $\Subst{\lambd[x][N]}{R}{y}$, kita akan mengandaikan
bahwa wakil $\lambd[x][N]$ dipilih sedemikian sehingga $x \neq y$ dan
$x \notin \FV{R}$.

Karena agak janggal menyebut $\lambd[x][x]$ sebuah ``kelas'', mulai sekarang
mari kita menyebut objek-objek tersebut suku-$\Lambda$ (atau cukup ``suku''
dalam bagian selanjutnya), untuk membedakannya dari suku-$\lambda$ yang telah
kita kenal.

\begin{editorial}
  Kita belum dapat meninggalkan suku: seluruh definisi suku-$\Lambda$
  didasarkan pada suku-$\lambda$, dan kita belum menyediakan metode untuk
  mendefinisikan fungsi pada suku-$\Lambda$. Artinya, semua fungsi semacam itu
  harus terlebih dahulu didefinisikan pada suku-$\lambda$, lalu
  ``diproyeksikan'' ke suku-$\Lambda$, sebagaimana kita lakukan untuk
  substitusi. Namun, kita mengandaikan pembaca dapat memahami secara intuitif
  bagaimana fungsi pada suku-$\Lambda$ dapat didefinisikan.
\end{editorial}
\end{document}
